% Part: lambda-calculus
% Chapter: church-rosser
% Section: parallel-beta-eta-reduction

\documentclass[../../../include/open-logic-section]{subfiles}

\begin{document}

\olfileid{lam}{cr}{pbe}

\olsection{समांतर $\beta\eta$-न्यूनीकरण}

या विभागात $\beta\eta$-न्यूनीकरणाशी निगडित समांतर
$\beta\eta$-न्यूनीकरणाला चर्च–रॉसर गुणधर्म आहे हे सिद्ध करू.

\begin{defn}[समांतर $\beta\eta$-न्यूनीकरण, $\beredpar$] \ollabel{defn:beredpar}
  पदांवरील \emph{समांतर $\beta\eta$-न्यूनीकरण}
  ($\beredpar$) पुढीलप्रमाणे विगमनाने परिभाषित केले जाते:
  \begin{enumerate}
    \item \ollabel{defn:beredpar1} $x \beredpar x$.
    \item \ollabel{defn:beredpar2} $N \beredpar N'$ असेल, तर $\lambd[x][N] \beredpar
      \lambd[x][N']$.
    \item \ollabel{defn:beredpar3} $P \beredpar P'$ आणि $Q \beredpar Q'$ असतील, तर $PQ \beredpar
      P'Q'$.
    \item \ollabel{defn:beredpar4} $N \beredpar N'$ आणि $Q \beredpar Q'$ असतील, तर
      $(\lambd[x][N])Q \beredpar \Subst{N'}{Q'}{x}$.
    \item \ollabel{defn:beredpar5} $N \beredpar N'$ असेल, तर $\lambd[x][Nx]
      \beredpar N'$, मात्र $x \notin FV(N)$ असले पाहिजे.
  \end{enumerate}
\end{defn}

\begin{thm}\ollabel{thm:refl}
  $M \beredpar M$.
\end{thm}

\begin{proof}
  सरावासाठी.
\end{proof}

\begin{prob}
  \olref[lam][cr][pbe]{thm:refl} सिद्ध करा.
\end{prob}

\begin{defn}[$\beta\eta$-पूर्ण विकास]\ollabel{defn:becd}
  $M$ चा \emph{$\beta\eta$-पूर्ण विकास}~$\becd{M}$
  पुढीलप्रमाणे परिभाषित केला जातो. पाचवी अट लागू
  होत असल्यास दुसऱ्या अटीऐवजी ती वापरायची आहे:
  \begin{align}
    \becd{x} &= x \ollabel{defn:becd1} \\
    \becd{(\lambd[x][N])} &= \lambd[x][\becd{N}] \ollabel{defn:becd2}\\
    \becd{(PQ)} &= \becd{P}\becd{Q} && \text{जर $P$ हे $\lambd$-अमूर्तीकरण नसेल तर}
    \ollabel{defn:becd3} \\
    \becd{((\lambd[x][N])Q)} &= \Subst{\becd{N}}{\becd{Q}}{x}
                             \ollabel{defn:becd4} \\
    \becd{(\lambd[x][Nx])} &= \becd{N} \ollabel{defn:becd5} & \text{जर $x
                                                      \notin FV(N)$}
  \end{align}
\end{defn}

\begin{lem}\ollabel{lem:comp}
  $M \beredpar M'$ आणि $R \beredpar R'$ असतील, तर $\Subst{M}{R}{y}
  \beredpar \Subst{M'}{R'}{y}$.
\end{lem}

\begin{proof}
  सर्व अमूर्तीकरणांची नावे आदेशनाच्या चरापासून आणि
  दोन्ही आदेशित पदांच्या मुक्त चरांपासून वेगळी निवडू.
  $M \beredpar M'$ च्या !!{derivation} वर विगमन करू.

  पहिली चार प्रकरणे \olref[pb]{lem:comp} मधील
  प्रकरणांसारखीच आहेत. अखेरचा नियम
  \olref{defn:beredpar5} असेल, तर काही $x$ आणि~$N'$
  साठी $M$ हे $\lambd[x][Nx]$, $M'$ हे $N'$
  आहे, $x \notin FV(N)$ आणि $N \beredpar N'$.
  आपल्याला $\Subst{(\lambd[x][Nx])}{R}{y} \beredpar
  \Subst{N'}{R'}{y}$, म्हणजे
  $\lambd[x][\Subst{N}{R}{y} x] \beredpar \Subst{N'}{R'}{y}$,
  हे सिद्ध करायचे आहे. ते \olref{defn:beredpar}\olref{defn:beredpar5}
  आणि विगमन गृहीतकावरून मिळते. बाह्य बद्ध नाव मूळ
  देहात मुक्त नाही आणि आदेशित पदातही मुक्त नाही, म्हणून
  आदेशनानंतरच्या देहात ते मुक्त होत नाही. त्यामुळे वापरलेल्या
  इटा नियमाची ताजेपणाची अट पूर्ण होते.
\end{proof}

\begin{lem}\ollabel{lem:cont}
  $M \beredpar M'$ असेल, तर $M' \beredpar \becd{M}$.
\end{lem}

\begin{proof}
%READERNOTE{SOL6-A126: इटा संकुचनाने फलनस्थानील अमूर्तीकरण नाहीसे होऊ शकते. त्यामुळे बीटा पुराव्यातील पहिली चार प्रकरणे तशीच लागू होत नाहीत. मुख्य triangle दावा आणि अमूर्तीकरणाच्या उपयोजनाचा सहाय्यक दावा यांचे एकत्र रचनात्मक विगमन देऊन ही उणीव पूर्ण केली आहे.}
प्रथम समांतर न्यूनीकरणाने नवीन मुक्त चरे निर्माण होत नाहीत,
हे त्याच्या पाच नियमांवरील विगमनाने मिळते. बीटा प्रकरणात
आदेशनाचे मुक्त-चर नियम वापरतात; इटा प्रकरणात बद्ध चर
उरलेल्या पदात मुक्त नसतो. पूर्ण विकास हाही एक समांतर
न्यूनीकरण आहे, हे त्याच्या व्याख्येवरील विगमनाने मिळते.
म्हणून पूर्ण विकासानेही मुक्त चरे वाढत नाहीत. खाली आवश्यक
तेथे प्रतिनिधींची बद्ध नावे सर्व संबंधित सांत मुक्त-चर
संचांपासून वेगळी निवडू.
पदाच्या रचनेवर पुढील दोन दावे एकत्र सिद्ध करू:
\begin{enumerate}
\item मुख्य दावा: $M\beredpar M'$ असल्यास
$M'\beredpar\becd{M}$.
\item सहाय्यक दावा: $A=\lambd[x][B]$, $A\beredpar A'$
आणि $S\beredpar T$ असल्यास
$A'S\beredpar\Subst{\becd{B}}{T}{x}$.
\end{enumerate}
सहाय्यक दावा अमूर्तीकरणांपुरताच आहे. तो प्रथम पाहू.
$A\beredpar A'$ ची शेवटची पायरी अमूर्तीकरणाचा नियम
असेल, तर $A'=\lambd[x][B']$ आणि $B\beredpar B'$.
लहान पदासाठीच्या मुख्य विगमन दाव्याने
$B'\beredpar\becd{B}$ मिळते. याला $S\beredpar T$
जोडून बीटा समांतर नियमाने सहाय्यक निष्कर्ष मिळतो.
शेवटची पायरी इटा नियमाची असेल, तर $B=Hx$,
$x\notin\FV{H}$, $A'=H'$ आणि $H\beredpar H'$.
$H$ अमूर्तीकरण नसेल, तर लहान पदासाठीच्या मुख्य दाव्याने
$H'\beredpar\becd{H}$ मिळते. उपयोजनाच्या नियमाने
$H'S\beredpar\becd{H}T$; ही उजवी बाजू
$\Subst{\becd{(Hx)}}{T}{x}$ आहे.
$H=\lambd[y][L]$ असेल, तर लहान अमूर्तीकरणासाठीचा
सहाय्यक विगमन दावा
$H'S\beredpar\Subst{\becd{L}}{T}{y}$ देतो.
येथे $y$ हे $x$ पासून आणि $T$ च्या मुक्त चरांपासून
वेगळे निवडले आहे. पूर्ण विकासाच्या चौथ्या नियमाने
$\becd{(Hx)}=\Subst{\becd{L}}{x}{y}$; तसेच $x$
हे $L$ मध्ये किंवा त्याच्या पूर्ण विकासात मुक्त नाही.
\olref[pb]{lem:comp} च्या सिद्धतेत दिलेल्या आदेशनसूत्राने
\[
\Subst{\Subst{\becd{L}}{x}{y}}{T}{x}
=\Subst{\becd{L}}{T}{y}.
\]
म्हणून या प्रकरणातही सहाय्यक दावा सिद्ध होतो.

आता मुख्य दावा सिद्ध करू. चरासाठी तो परावर्तितेने मिळतो.
$M=PQ$ अमूर्तीकरण नसलेल्या फलनस्थानील पदासह असेल,
तर समांतर उपयोजनाचाच नियम लागू होतो. दोन्ही उपपदांना
मुख्य विगमन दावा लावून $P'Q'\beredpar\becd{P}\becd{Q}$
मिळते; उजवी बाजू $\becd{M}$ आहे.
$M=(\lambd[x][N])Q$ असेल आणि शेवटची पायरी बीटा
नियमाची असेल, तर $M'=\Subst{N'}{Q'}{x}$,
$N\beredpar N'$ आणि $Q\beredpar Q'$.
दोन्ही मुख्य विगमन दावे आणि \olref{lem:comp} यांनी
$M'\beredpar\Subst{\becd{N}}{\becd{Q}}{x}=\becd{M}$
मिळते. शेवटची पायरी उपयोजनाची असेल, तर
$M'=P'Q'$, $\lambd[x][N]\beredpar P'$ आणि
$Q\beredpar Q'$. मुख्य विगमन दाव्याने
$Q'\beredpar\becd{Q}$. मग लहान फलनस्थानील
अमूर्तीकरणासाठीचा सहाय्यक दावा तोच आवश्यक निष्कर्ष देतो;
यात $P'$ अजून अमूर्तीकरण असणे गृहीत धरलेले नाही.
शेवटी $M=\lambd[x][N]$ पाहू. शेवटची पायरी इटा
नियमाची असेल, तर $N=Hx$, $x\notin\FV{H}$ आणि
$M'=H'$ जेथे $H\beredpar H'$. मुख्य विगमन दाव्याने
$H'\beredpar\becd{H}=\becd{M}$.
शेवटची पायरी अमूर्तीकरणाच्या नियमाची असेल, तर
$M'=\lambd[x][N']$ जेथे $N\beredpar N'$.
$M$ इटा रेडेक्स नसेल, तर मुख्य विगमन दावा आणि
अमूर्तीकरणाचा नियम थेट निष्कर्ष देतात.
इटा रेडेक्स असेल, तर $N=Hx$ आणि $x\notin\FV{H}$.
$Hx\beredpar N'$ ची शेवटची पायरी उपयोजनाची असेल,
तर $N'=H'x$, $H\beredpar H'$; मुख्य विगमन दाव्याने
$H'\beredpar\becd{H}$ आणि इटा नियमाने
$\lambd[x][H'x]\beredpar\becd{H}=\becd{M}$.
ही पायरी बीटाची असेल, तर $H=\lambd[y][L]$,
$L\beredpar L'$ आणि $N'=\Subst{L'}{x}{y}$.
नवीन मुक्त चरे निर्माण होत नसल्याने $x\notin\FV{L'}$,
आणि $x\ne y$ अशी नावांची निवड आहे. म्हणून
$\lambd[x][\Subst{L'}{x}{y}]$ हे $\lambd[y][L']$
याच वर्गपदाचे दुसरे रूप आहे. अमूर्तीकरणाच्या नियमाने
$H\beredpar\lambd[y][L']$, आणि लहान पद $H$ साठीचा
मुख्य विगमन दावा
$\lambd[y][L']\beredpar\becd{H}=\becd{M}$ देतो.
ही सर्व शक्य प्रकरणे आहेत. प्रत्येक वापरलेल्या विगमन
दाव्याचे मूलपद सध्याच्या पदाचे योग्य उपपद आहे,
म्हणून एकत्र विगमन पूर्ण झाले.
\end{proof}

\begin{thm}\ollabel{thm:cr}
  $\beredpar$ ला चर्च–रॉसर गुणधर्म आहे.
\end{thm}

\begin{proof}
  \olref{lem:cont} वरून लगेच मिळते.
\end{proof}

\end{document}
